\documentclass[a4paper]{report}

%====================== PACKAGES ======================

\usepackage[french]{babel}
\usepackage[utf8]{inputenc}
%pour gérer les positionnement d'images
\usepackage{float}
\usepackage{amsmath}
\usepackage{graphicx}
\usepackage[colorinlistoftodos]{todonotes}
\usepackage{url}
%pour les informations sur un document compilé en PDF et les liens externes / internes
\usepackage{hyperref}
%pour la mise en page des tableaux
\usepackage{array}
\usepackage{tabularx}
%pour utiliser \floatbarrier
%\usepackage{placeins}
%\usepackage{floatrow}
%espacement entre les lignes
\usepackage{setspace}
%modifier la mise en page de l'abstract
\usepackage{abstract}
%police et mise en page (marges) du document
\usepackage[T1]{fontenc}
\usepackage[top=2cm, bottom=2cm, left=2cm, right=2cm]{geometry}
%Pour les galerie d'images
\usepackage{subfig}

\usepackage{blindtext}


%====================== INFORMATION ET REGLES ======================

%rajouter les numérotation pour les \paragraphe et \subparagraphe
\setcounter{secnumdepth}{4}
\setcounter{tocdepth}{4}

\hypersetup{							% Information sur le document
pdfauthor = {Nicolas Renard},				% Auteurs
pdftitle = {Nom du Projet -
			Sujet du Projet},			% Titre du document
pdfsubject = {Mémoire de master},		% Sujet
pdfkeywords = {securite,alienvault,atheo},	% Mots-clefs
pdfstartview={FitH}}					% ajuste la page à la largueur de l'écran
%pdfcreator = {MikTeX},% Logiciel qui a crée le document
%pdfproducer = {}} % Société avec produit le logiciel

%======================== DEBUT DU DOCUMENT ========================

\begin{document}

%régler l'espacement entre les lignes
\newcommand{\HRule}{\rule{\linewidth}{0.5mm}}

%page de garde
\begin{titlepage}
\begin{center}

% Upper part of the page. The '~' is needed because only works if a paragraph has started.
\includegraphics[width=0.35\textwidth]{../img/logo-uds}~\\[0.5cm]
\includegraphics[width=0.35\textwidth]{../img/logo_menu.png}~\\[1cm]

\textsc{\LARGE Université de Strasbourg} \\[0.5cm]
\textsc{\LARGE Athéo Ingénierie}\\[1.5cm]


\textsc{\Large }\\[0.5cm]

% Title
\HRule \\[0.4cm]

{\huge {\bfseries Mémoire de master Réseaux Informatiques et Systèmes Embarqués.}\\
Analyse, mise en place et mise en vente d'une solution de sécurité unifiée. \\[0.4cm] }

\HRule \\[1.5cm]

% Author and supervisor
\begin{minipage}{0.4\textwidth}
\begin{flushleft} \large
\emph{Auteur:}\\
Nicolas \textsc{Renard}\\
\end{flushleft}
\end{minipage}
\begin{minipage}{0.4\textwidth}
\begin{flushright} \large
\emph{Client:} \\
Benjamin \textsc{His}\\
\emph{Référent:} \\
Julien \textsc{Montavont}
\end{flushright}
\end{minipage}

\vfill

% Bottom of the page
{\large \today}

\end{center}
\end{titlepage}

%page blanche
\newpage
~
%ne pas numéroter cette page
\thispagestyle{empty}
\newpage

\renewcommand{\abstractnamefont}{\normalfont\Large\bfseries}
%\renewcommand{\abstracttextfont}{\normalfont\Huge}

\begin{abstract}
\hskip7mm

\begin{spacing}{1.3}

Ce rapport de stage clôture deux années de master Résaux Informatiques et Systèmes embarqués à l'Université de Strasbourg. Dans celui-ci, je vais présenter et développer le travail réalisé durant ce stage. Ma mission consistait en l'analyse, la mise en place et la mise en vente d'une solution de sécurité unifiée. M'intéressant à la sécurité depuis très longtemps, j'ai été très intéressé par ce sujet dès le début du stage. \\
Je vais diviser ce rapport en plusieurs parties. Tout d'abord je présenterai le contexte dans lequel le stage a été réalisé, en présentant les activités de l'entreprise et plus précisément les personnes avec lesquelles j'ai travaillé. Ensuite j'expliquerai le contexte global de la sécurité dans le domaine informatique et j'en démontrerai son importance. Puis j'effectuerai un comparatif de différents produits et je présenterai le produit sur lequel j'ai travaillé, son architecture et ses fonctionnalités. Finalement je conclurai ce rapport en proposant de nouvelles perspectives concernant la sécurité et en expliquant ce que ce stage m'a apporté.

\end{spacing}
\end{abstract}


\tableofcontents
\thispagestyle{empty}
\setcounter{page}{0}
%ne pas numéroter le sommaire

\newpage

%espacement entre les lignes d'un tableau
\renewcommand{\arraystretch}{1.5}

%====================== INCLUSION DES PARTIES ======================

~
\thispagestyle{empty}
%recommencer la numérotation des pages à "1"
\setcounter{page}{0}
\newpage

\chapter{Introduction}
Dans ce chapitre, je présenterai la société Athéo ainsi que ses activités. Je définirai également le contexte dans lequel j'ai réalisé ma mission de stage. 

\section{Présentation de l'entreprise}
Athéo\footnote{http://www.atheo.net} est une entreprise de services du numérique basée à Strasbourg. Elle est composée d'une vingtaine de personnes, principalement des ingénieurs informatique. Athéo est spécialisée dans l'intégration de systèmes et de réseaux informatiques, l'infogérance, le monitoring et la sécurité. Une grande partie de ses collaborateurs est issue de la liquidation de la société de services RBS. Ainsi, dès son lancement, Athéo celle-ci a pu compter sur les compétences d'experts aguerris. \\

Lors de mon arrivée dans la société, j'ai été accueilli par mon tuteur de stage, Benjamin \textsc{His}, directeur technique et architecte système. En plus de diriger l'équipe d'ingénieurs, il s'occupe également de toutes les offres de sécurité d'Athéo. Aux cotés de Benjamin, j'ai également travaillé avec deux experts réseau, Erkan \textsc{Valentin} et Stéphane \textsc{Merl}. 

\section{Mission de stage}
Cette partie présente la mission principale ainsi que les tâches secondaires réalisées durant le stage. \\
Ma mission principale consistait en l'analyse, la mise en place et la mise en vente d'une solution de sécurité unifiée. Ce programme est un gestionnaire d'informations et d'évènements de sécurité (SIEM\footnote{Security Information and Event management} en anglais) appelé Alienvault\footnote{https://www.alienvault.com/}. Ce programme installé sur un serveur surveille tous les équipements présents dans le réseau et détecte les incidents de sécurité, affiche des statistiques et génère des rapports à partir des données collectées. Les fonctionnalités apportées par Alienvault seront discutées plus en profondeur dans la partie \ref{presentationalienvault}.\\
En plus de ma mission principale, j'ai effectué d'autres tâches secondaires durant mon stage. 

\section{Tâches secondaires réalisées}
% TODO %
% Description des autres tâches que j'ai faites:
% - configuration de switch
% - RDV client
% - cisco prime
% - phishing
\blindtext\\



\chapter{Contexte}
%TODO%
% - Pourquoi la sécurité informatique?
% - Coûts de perte de données
% - cf. étude Institut Ponemon


\blindtext\\

\blindtext\\

\blindtext\\
 
\chapter{Pourquoi la sécurité}
\blindtext\\

\section{Alienvault}
\label{presentationalienvault}

\blindtext\\

\blindtext\\

\blindtext\\

\chapter{Autre partie}


\blindtext\\

\blindtext\\

\blindtext\\

\chapter{Résultats}

\blindtext\\

\blindtext\\

\blindtext\\

\chapter{Bilan}

%Rappel du context
Intro / Rappel Contexte \\

\blindtext\\

\blindtext\\


%Ne pas numéroter cette partie
\part*{Annexes}
%Rajouter la ligne "Annexes" dans le sommaire
\addcontentsline{toc}{part}{Annexes}

\chapter*{Annexe 1}
\addcontentsline{toc}{chapter}{Annexe 1}

%changer le format des sections, subsections pour apparaittre sans le num de chapitre
\makeatletter
\renewcommand{\thesection}{\@arabic\c@section}
\makeatother

\blindtext[1]\\

\blindlistlist[2]{itemize}[6]

\blindtext[1]

\chapter*{Annexe 2}
\addcontentsline{toc}{chapter}{Annexe 2}

%recommencer la numérotation des section à "1"
\setcounter{section}{0}

\blindtext[1]\\

\blindlist{itemize}[6]

\blindtext[1]\\

\newpage

%récupérer les citation avec "/footnotemark"
\nocite{*}

%choix du style de la biblio
\bibliographystyle{plain}
%inclusion de la biblio
\bibliography{bibliographie.bib}
%voir wiki pour plus d'information sur la syntaxe des entrées d'une bibliographie

\end{document}